\documentclass[ms,a4paper]{memoir}
\chapterstyle{dash} % try also reparticle
\usepackage{ulem}   % underline
\usepackage{xcolor}
\usepackage{graphicx}
\graphicspath{ {./img/} }
%% Polski
\usepackage[T1]{fontenc}
\usepackage[polish]{babel}
\usepackage[utf8]{inputenc}
\usepackage{marginnote}
%
%\usepackage{hyperref}
% All font size must be normal size
\renewcommand{\large}{\normalsize}
\renewcommand{\Large}{\normalsize}
\newcommand{\red}[1]{\textcolor{red!50!black}{#1}}
\newcommand{\RED}[1]{\textcolor{red!50!black}{\MakeUppercase{#1}}}
% font hyphenation
\usepackage{everysel}
\EverySelectfont{%
\fontdimen2\font=0.6em % interword space
\fontdimen3\font=0.2em % interword stretch
\fontdimen4\font=0.1em % interword shrink
\fontdimen7\font=0.9em % extra space
\hyphenchar\font=`\-% to allow hyphenation
}


%\usepackage[spacing=true,factor=1200, stretch=10, shrink=15]{microtype}

\renewcommand{\baselinestretch}{1}


\newbox\flinebox{}
\newbox\slinebox{}
\newbox\mlinebox{}
\def\duplines{\setlength\parindent{0pt}
  \setbox\flinebox\lastbox{}
  \ifvoid\flinebox\relax
  \else
  \setbox\slinebox\hbox{\copy\flinebox}
  \setbox\mlinebox\hbox{\copy\flinebox}
  \unskip\unpenalty{}
  {\duplines}

{\color{black!30} \box\flinebox\vspace*{-2.85ex}}
{\color{black!50} \makebox[\textwidth]{\hspace*{-0.25pt}\box\mlinebox}\vspace*{-2.75ex}}
{\color{black!90}  \makebox[\textwidth]{\hspace*{0.25pt}\box\slinebox}}\fi

}

\newcommand\BlurText[1]{%
  \vbox{#1\par\duplines}}


\begin{document}
\chapter{\RED{Skrót}}
\section{}
\paragraph{}
Pakaczuki jest przygodą pod realia Zewu Cthulhu, rozgrywać się powinna w międzywojniu (1920 \- 1930)
w centralnych Stanach Zjednoczonych. Przewiduje się ją na \emph{1} do 2 sesji, przy ekipie od \emph{1} do 4 graczy.\\
\section{Skrót}
Gracze na zaproszenie wuja jadą do małej miejscowości w górach (Pakaczuki).
Droga do tego miejsca jest utrudniona z powodu na rozmokłe drogi oraz pozawalane mosty.
Na miejscu okazuje się, że wiosenne roztopy co roku zawalają jedyny most z miejscowości, a mieszkańcy składają przybyłych w ofierze.
Jedyne możliwe wyjście z terenu osady jest przez kopalnie i  nieeuklidesowy labirynt.
\paragraph{}
Przygoda jest aktualnie w stanie \red{alpha}.
\newpage
\tableofcontents
\newpage

\chapter{\RED{Przed rozgrywką}}

\section{Wymagania co to graczy}
Przygoda nie ma jakiś szczególnych wymagań co do graczy.
Wszelkie niezbędne przedmioty będą do znalezienia na miejscu,
a żadna z umiejętności nie jest niezbędna do pchnięcia przygody
na przód. Jedynym ważnym elementem jest fakt, że jeden z graczy
ma postać nazywającą się \emph{Evert Hampton}. Wymóg ten jest
podyktowanym treścią testamentu.
\section{Potrzebne materiały}
Mistrz gry powinien wyposażyć się w:
\begin{itemize}
\item Małą kopertę
\item Dużą kopertę (w której zmieści się mała)
\item Wydruk testamentu ze strony %\pageref{hand:testament}
\end{itemize}

\chapter{\RED{Prolog}}
\section{Tło fabularne}
Jeden z graczy (Evert Hampton) otrzymał \emph{tydzień temu} list od swojego wuja, Johna.
Wuj zaprasza go do Pakaczuki, swojego domu, na obchody \emph{Święta}.
Jest to pierwszy list od roku i nie ma w nim wzmianki o przepisie na pierogi, o który Evert prosił.
Nie jest to jednak nic dziwnego, wuj w przeszłości potrafił przestać pisać na parę miesięcy i po tym czasie odpisać jak by nic się nie stało.
Podekscytowany możliwością spotkania krewnego, którego nie widział od ponad
siedmiu lat, Evert chętnie zabiera się do drogi.
Wuj w liście zaprasza nie tylko jego, ale jeśli ma jeszcze z dwóch trzech przyjaciół, to spokojnie się zmieścicie.


\section{Otwarcie}
Postaci podróżują już jeden lub dwa dni do celu.
Podróż przebiega w miarę normalnie. Ford T którym się przemieszczają nie zawodzi, a pogoda
dopisuje. Mimo że jest \emph{przełom marca i lutego} słońce dogrzewa i jedzie się stosunkowo
przyjemnie. Właśnie dojechali do zajazdu w małej
miejscowości pod górami. [TODO] bo tak nazywa się ta miejscowość. Jest praktycznie ostatnim miejscem z ogólnie rozumianą cywilizacją.
Miejscowość ta ma poza zajazdem kilka domów, sklep i warsztat/stację benzynową.
Zwróć uwagę graczy, że wedle mapy stąd do Wuja jest jeszcze 2h.
Jak wezmą poprawkę na pogodę i średnie warunki drogowe, maksymalnie 4h jazdy.
Trzeba zrobić wrażenie, że nie da się już skoczyć tam teraz po nocy, ale rano to moment i są na miejscu.

\subsection{Miejsce: Zajazd}
Zajazd jest wielkości dużego domu.
Sprawia wrażenie miejsca raczej przytulnego.
Ma dwa piętra i parter.
Na piętrach jest 8 pokoi do spania.
A parter przeznaczony jest na jadalnie, recepcję i pokój oczekiwania.
Z uwagi na fakt, że gracze przyjeżdżają totalnie poza jakimkolwiek sezonem, są jednymi gośćmi.
Miejsce to prowadzi rodzina mieszkająca parę domów dalej, akurat jak gracze przyjadą na recepcji
jest Pani Hanna Sevani.
\marginnote{|Osoba: Hanna Sevani}
Hanna jest starszą ale jeszcze aktywną kobietą, pod 50 lat.
Jest aż zbyt słodka, może odnieść się do postaci per `kochanieńki' jeśli jej pozwolimy.
Pyta co gracze chcą na śniadanie, ale z uwagi na prywatne powody miękko upiera się by to było o 7 rano nie później.
\\
Tu mamy scenę śniadania w zajeździe.
Jest to dobre miejsce na opisanie wyglądu postaci i krótką pogawędkę przy jajecznicy.

Przy samym wyjeździe zwróć uwagę, że piękny wschód słońca przeszedł w równie piękny poranek.
Ptaszki śpiewają a rzadkie przebiśniegi i pęki drzew wydają się im zalewać cały świat zielonością.
Pilnuj by nie zabrali więcej niż jednej kanapki na drogę, może wspomnij że Wujo pewnie ich ugości.

Wyjechali?
\chapter{\RED{Droga do Pakaczuki}}

Jak tylko zajazd zniknie im za tylną szybą na tyle by nie myśleli wracać, zasnuj niebo szarymi chmurami.
Jest bardzo ciężka pogoda, wieje, mży, a koła Forda wprawionego w takich warunkach ślizgają się na błocie pośniegowym.
Turlają się tak z 6 do 8 godzin.
Opisz jak bardzo jest im zimno.
Jak bardzo czują że ich kości zamarzają.
Jak przeklinają tą przemoczoną mapę na której już nic nie widać.
Jeśli zarzucą, że mają takie umiejętności jak ``Nawigacja'' wyjaśnij im, że nie robiłeś testów bo nawigacja idzie im świetnie.
Tylko po drodze mijają już któryś zawalony most i próbują na około.
Nie zatrzymali się też po drodze, ugotować coś albo zrobić miejsce do spania bo za każdym razem wierzyli, że to już za tym wzgórzem bo mapa pokazuje.
Jak już osłabiłeś ich ducha.
Jak już faktycznie czują beznadzieje sytuacji czas wprowadzić kolejnego BNa.
\paragraph{}
\marginnote{|Osoba:Wolny Jeleń}
Z kierunku w którym jadą nadjeżdża na koniu Indianin.
Jest pozytywnie nastawiony i macha do graczy.
Odziany jest w Jeansy i Ponczo.
Gdy podjedzie zauważa, że pewnie się zgubili i dalsza jazda tą drogą nie ma najmniejszego sensu, tamten most też jest zawalony.
Po usłyszeniu, że jadą do Pakaczuki, tłumaczy że zna drogę, ale teraz to by się tam pieszo by nie pakował, a automobilem to już na pewno.
Zauważa też, że idzie burza, proponuje w geście dobrej woli, że ich przenocuje i nalega by iść za nim.
Może poprosić o podwiezienie ``futer'', które udało mu się upolować na bagażniku samochodu.
Spokojnie ma koc którym wyłożył grzbiet konia więc nie ubrudzi.
Musisz na graczach wywrzeć presję by z nim poszli.
Jeśli nie będą chcieli z nim pójść daj im pobłądzić jeszcze chwilę i znaleźć jego chatę.
\marginnote{|Miejsce: Chata wolnego Jelenia}
Jest to duża, jednoizbowa drewniana chata w stylu ``traperskim''.
Miejsce jest bardzo przytulne, zapach gulaszu z dziczyzny i ciepło kominka poprawia nastrój.
Jeleń wchodząc wita serdecznie żonę i dwoje dzieciaków.
Prosi ją po indiańsku o coś, a postaciom graczy poleca zasiąść za dość sporym stołem.
Po chwili Pan domu kładzie na stole dzban herbaty, to jest ta chwila gdy robisz to samo jako mistrz gry.
Ten mały trik ma spowodować by faktycznie poczuli się pewnie.
Gdy gracze rozmarzają przy herbatce Jeleń przeprasza ich, znika na zewnątrz i po chwili wraca ze zdobyczami.
Siada w kącie chaty i zaczyna rozmowę z postaciami graczy, skórując króliki.
Na stół trafia gulasz, po dobrej gorącej porcji dla każdego.
Jeleń płynnie przechodzi do historii o Pakaczuki.
\marginnote{|Historia Pakaczuki}
\fbox{ \begin{minipage}{32em}
    Pakaczuki jest to miejscowość, założona przez białych w XVIII.
    Wcześniej była tam ``święta ziemia'' Indian.
    Głównie dlatego że nic sensownego tam nie dało się robić.
    Zwierzyny mało, las taki sobie.
    Jedyne co było tam używane to stos pogrzebowy.
    I z tą świętą ziemią też nie należy przesadzać.
    Nazwa nadana przez białych pochodzi od tego jak w żartach nazywali ją Indianie.
    Paka znaczyło u nich czoło a Czuki płaskie.
    Wedle legendy ówczesny wódz plemienia sprzedał je przyjezdnym.
    Tak, sprzedał.
    Znaczyło to tyle, że biali, co było niespotykane, zaproponowali za ten kawałek ziemi, wiele cennych skór, kilka grzmiących kijów itd.
    Wódz zgodził się, ale tak go rozbawiła i zdziwiła ta propozycja, że uderzył się w czoło otwartą dłonią.
    Stąd nazwa.
    Jedynym przeciwnym był stary szaman mówiący coś o przedwiecznych tunelach których należy strzec.
    Ale nikt go nie słuchał.
\end{minipage}}
\paragraph{}
Po tej historii, zwraca uwagę, że w sumie jest już późno.
A rano rusza na polowanie.
Chętnie pokaże graczom jak dostać się do wspomnianej miejscowości, bo czemu by w tamtej okolicy nie wybrać się jutro na polowanie.
Pokazuje im miejsce do spania, które jest naprawdę sporą stertą świeżych futer, ale nie wydobywa się z tego przykry zapach.
\paragraph{}
Wstają wypoczęci, na chwilę przed świtem.
Droga zajmie im nie długo i około 8 powinni być na moście wjazdowym.
Tam też Jeleń się żegna i idzie w swoją stronę.
Jeśli uważasz, że mało było turlania koścmi można w tym miejscu zrobić jakiś test prowadzenia.
Jednak karą za niezdanie powinna być mask ryski na lakierze czy śmiech przewodnika.


\chapter{\RED{Pakaczuki}}

\paragraph{}
Pakaczuki jest miejscowością położoną w dolince otoczonej z każdej strony stromymi ścianami skalnymi.
Do opisu można spokojnie użyć słowa miska, czy nawet szklanka.
Jeśli jakiś gracz ma cokolwiek w umiejętności ``wspinaczka'' zwróć uwagę, że nie pakował by się na te ściany teraz, bo była by to pewna śmierć.
Miejscowi zdają się patrzeć podejrzliwie, ale nie jest to szczególnie inne od innych zabitych dechami dziur.
\section{Miejsca}
\subsection{mapa}
\begin{figure}
  \includegraphics[width=\linewidth]{mapa}
  \end{figure}

\subsection{Most}
\paragraph{}
Most jest chyba najważniejszym elementem wioski.
Jest on szeroki i drewniany, wykonany solidnie.
Jednak wyraźnie widać, że nie ma on więcej niż rok.
Drewno nie jest malowane ani olejowane, jak by nie miał stać tam dość długo by zasłużyć na takie zapieczenie.
Rzeka pod nim jest dość rwąca i szeroka, przepływanie wpład odpada.
\paragraph{}
Zawalenie się mostu jest ważnym wydażeniem w życiu miejscowości.
Powinno ono nastąpic jak tylko gracze uznają, że chcą wracać, ale nie później niż drugiego dnia.
\marginnote{|Święto}
Gdy tylko to nastąpi przez Pakaczuki przebiega młody chłopiec krzyczący, że most się zawalił.

\subsection{cmentarz}



%\label{hand:testament}
\end{document}
